\section{Diseño preliminar de Sistemas de Control I}

\subsection{Variable controlada (PV)}

La variable controlada es la temperatura del fluido de proceso (agua o agua con detergente diluido) en el tanque de mezcla del sistema CIP portátil, medida mediante un sensor RTD PT100 conectado a un módulo de expansión analógica SM 1231 (escalamiento \texttt{NORM\_X/SCALE\_X} en TIA Portal).

De las variables de proceso disponibles (temperatura, caudal, nivel, presión y conductividad), la temperatura es la única con una dinámica lo suficientemente lenta y bien caracterizable con los recursos del laboratorio para diseñar y validar un controlador P/PI/PID completo dentro del alcance del curso de Sistemas de Control I.

\subsection{Definición del lazo de control}
\begin{table}[H]\centering\small\caption{Definición preliminar del lazo de control.}
\begin{tabularx}{\textwidth}{L{3.7cm}YL{3cm}}\toprule
\textbf{Elemento}&\textbf{Definición}&\textbf{Unidad/rango}\\\midrule
Variable controlada&Temperatura del agua o solución&\campo{$^{\circ}$C y rango}\\
Variable manipulada&Potencia entregada al calentador&\campo{W o \%}\\
Perturbaciones&\campo{Pérdidas, caudal, temperatura de entrada}&\campo{Unidades}\\
Punto de operación&\campo{Valor nominal}&\campo{Unidad}\\
Limitaciones&\campo{Saturación y límites de seguridad}&\campo{Valores}\\\bottomrule
\end{tabularx}\end{table}


\subsection{Diagrama de bloques}
\insertarfigura{diagrama_bloques_control.pdf}{Diagrama de bloques del sistema térmico.}{fig:bloquescontrol}

\subsection{Modelo dinámico preliminar}
\campo{Presente el balance de energía, parámetros, condiciones iniciales, función de transferencia o modelo equivalente, punto de operación y rango de validez.}

\begin{equation}
\boxed{\text{Inserte aquí la ecuación principal del modelo}}
\end{equation}

\subsection{Simulación en lazo abierto}
\insertarfigura{respuesta_lazo_abierto.pdf}{Respuesta preliminar del modelo en lazo abierto.}{fig:lazoabierto}
\campo{Analice ganancia, constante de tiempo, retardo, estabilidad y limitaciones observadas.}

\subsection{Especificaciones preliminares de desempeño}
\begin{table}[H]\centering\small\caption{Especificaciones preliminares de desempeño.}
\begin{tabularx}{\textwidth}{YL{3cm}L{4cm}}\toprule
\textbf{Indicador}&\textbf{Meta}&\textbf{Justificación o fuente}\\\midrule
Error estacionario&\campo{\%}&\campo{--}\\
Sobreimpulso&\campo{\%}&\campo{--}\\
Tiempo de establecimiento&\campo{s o min}&\campo{--}\\
Esfuerzo de control&\campo{Límite}&\campo{--}\\\bottomrule
\end{tabularx}\end{table}

